\begin{textsample}{POS Dim 2 – 2013\_c6 – Score -2.00 – 2013\_c6/enade\_2013\_fisio\_a31223fe5a.txt}  \label{ex:f2_pos_378}
CAPÍTULO 3 ANÁLISE TÉCNICA DA PROVA Este capítulo tem por objetivo apresentar o desempenho dos estudantes concluintes de Fisioterapia no ENADE/2013.
Para isso, foram calculadas as estatísticas básicas da prova em seu todo ( seção 3.1.1 ), bem como as estatísticas dos componentes relacionadas ao Componente de Formação Geral ( seção 3.1.2 ) e ao de Conhecimento Específico da Área ( seção 3.1.3 ).
Dadas as suas características, foram analisadas em separado, as questões objetivas ( seção 3.2 ) e as questões discursivas ( seção 3.3 ).
Para as questões objetivas, foram disponibilizados os índices de facilidade e o de discriminação bisserial, também em separado para os Componentes de Formação Geral ( seção 3.2.1 ) e de Conhecimento Específico ( seção 3.2.2 ).
De cada componente, uma das questões foi escolhida para exemplificar a análise gráfica relacionando as alternativas escolhidas pelos estudantes ( inclusive o gabarito ) com o número de acertos no componente.
O Anexo I apresenta a íntegra da análise gráfica para todas as questões objetivas.
Para cada uma das questões discursivas, os conteúdos dos tipos mais comuns de respostas dos alunos são apresentados e comparados com o padrão de resposta esperado ( ver Anexo VIII com o padrão de respostas ).
Tomando como base as duas questões discursivas do Componente de Formação Geral, a seção 3.3.1.6 apresenta comentários sobre a correção das respostas com respeito à Língua Portuguesa.
Nas tabelas, são apresentados o tamanho da população inscrita e de presentes, e as seguintes estatísticas das notas11 : média do desempenho na prova, erro padrão da média, desvio padrão, nota mínima, mediana e nota máxima.
Tais estatísticas contemplam o total de estudantes concluintes da Área de Fisioterapia inscritos e presentes à prova do ENADE/2013, tendo -se em vista as seguintes agregações : ( a ) as Grandes Regiões e o país como um todo ; ( b ) a Categoria Administrativa ; e ( c ) a Organização Acadêmica.
Em relação aos gráficos de distribuição de notas, o intervalo considerado foi de 10 unidades, aberto à esquerda e fechado à direita, com exceção do primeiro intervalo, [ 0 ; 10 ], fechado em ambos os extremos.
Para os gráficos de distribuição das notas das questões discursivas, foram consideradas mais duas categorias : questão em branco e nota zero. 11 Essas estatísticas e outras estão definidas no Capítulo 1. 3.1 ESTATÍSTICAS BÁSICAS DA PROVA 3.1.1 Estatísticas Básicas Gerais A Tabela 3.1 apresenta as estatísticas básicas da prova por grande Região.
A população total de inscritos foi de 12.420.
Destes, 11.106 estiveram presentes, sendo 10,6 \% o índice de não comparecimento.
A Região de maior abstenção foi a Norte ( 19,5 \% ) e a de menor abstenção foi a Sul ( 3,8 \% ).
A média das notas da prova como um todo ( nas seções seguintes serão analisados os Componentes de Formação Geral e de Conhecimento Específico ) foi 48,5, sendo que os alunos da região Norte obtiveram a média mais baixa ( 44,8 ) e os da região Sul obtiveram a média mais alta ( 51,3 ).
As demais médias foram : 47,1 na região Nordeste, 49,1 na região Sudeste e 47,3 na região Centro-Oeste.
O desvio padrão para o Brasil como um todo foi 11,5, sendo o maior desvio padrão encontrado na região Sudeste ( 11,7 ) e o menor na região Sul ( 10,7 ), indicando uma menor dispersão das notas desta última região.
A região que obteve a maior nota máxima foi a Sul ( 83,6 ), ao passo que a região que atingiu a menor nota máxima foi a Norte ( 76,0 ).
A mediana do Brasil como um todo foi 48,7, sendo a maior mediana obtida na região Sul ( 51,6 ) e a menor obtida na Norte ( 45,2 ).
A menor nota mínima foi zero nas regiões Norte e Nordeste, ao passo que nas demais regiões as notas mínimas foram : 9,3 na região Sudeste, 12,1 na região Sul e 7,7 na região Centro-Oeste.
O comportamento das notas dos estudantes de todo o Brasil pode ser observado no Gráfico 3.1 que apresenta um histograma com a distribuição das mesmas.
Essa é uma distribuição unimodal com moda no intervalo ( 40 ; 50 ].
Destaca -se, também, o intervalo ( 50 ; 60 ] com frequência relativa muito próxima à modal.
O coeficiente de assimetria da distribuição das notas é negativo ( -0,11 ), porém de pouca magnitude, e podemos considerar que a distribuição é aproximadamente simétrica.
As distribuições por Grande Região também apresentam assimetria negativa e de pouca magnitude, concentração pouco maior do lado direito do histograma e mais espalhada do lado esquerdo, excetuando -se apenas a região Centro-Oeste que possui assimetria positiva e também de pouca magnitude ( 0,10 ).
Os Gráficos 3.2, 3.3 e 3.4 apresentam informações referentes à média da nota final dos participantes, desagregando os resultados de acordo com, respectivamente, as Grandes Regiões do país, a Categoria Administrativa e a Organização Acadêmica.
Os gráficos apresentam o valor da média das notas como uma barra e os extremos do intervalo de confiança de 95 \% como linhas verticais unidas por uma linha horizontal na forma da letra H maiúscula.
Considerando -se o gráfico de notas segundo Grande Região ( Gráfico 3.2 ), observa -se que existe diferença estatisticamente significativa ao nível de 95 \% entre a maior média, obtida na região Sul ( 51,3 ) e a menor, obtida na região Norte ( 44,8 ).
Levando -se em conta os estudantes por Categorias Administrativas ( Gráfico 3.3 ), observa -se que existe diferença estatisticamente significante entre as médias das notas das IES Públicas e Privadas.
Pode ser observado ainda que a média dos alunos de IES Públicas ( 56,3 ) é maior que a dos alunos de IES Privadas ( 47,3 ).
A diferença entre as médias das regiões Sul e Norte ( 6,5 ), a maior e a menor média, é inferior a diferença entre IES Públicas e Privadas ( 9,0 ).
Tendo como foco o Gráfico 3.4, que apresenta as notas médias das provas segundo Organização Acadêmica, constata -se que a diferença da média das notas dos estudantes provenientes de Universidades é estatisticamente significativa ao nível de 95 \% das médias das notas dos estudantes provenientes de Centros Universitários e Faculdades, assim como a diferença entre a média dos estudantes de Centros Universitários e Faculdades.
A maior média foi obtida pelos estudantes de Universidades ( 50,7 ), e a menor, pelos de Faculdades ( 45,8 ). 3.1.2 Estatísticas Básicas no Componente de Formação Geral A Tabela 3.2 apresenta as estatísticas básicas em relação ao componente da prova que avalia a Formação Geral dos estudantes concluintes.
Os alunos de todo Brasil obtiveram desempenho médio de 45,8.
Quanto à variabilidade, o desvio padrão das notas dos estudantes do Brasil como um todo foi 14,4.
A maior média foi obtida na região Sudeste ( 46,7 ), e a menor, na região Norte ( 41,7 ).
As demais médias foram : 45,1 na região Nordeste, 46,2 na região Sul e 45,9 na região Centro-Oeste.
Já o maior desvio padrão foi obtido na região Nordeste ( 14,9 ) e o menor na região Centro-Oeste ( 13,7 ).
Os demais desvios padrões foram : 14,1 na região Norte, 14,2 na região Sudeste e 14,0 na região Sul.
A maior nota no Componente de Formação Geral da prova do ENADE foi obtida por pelo menos um aluno da região Centro-Oeste ( 88,4 ) enquanto que a menor nota máxima foi obtida na região Norte ( 80,3 ).
Nas outras regiões as notas máximas foram : 87,6 na região Nordeste, 88,0 na região Sudeste e 86,5 na região Sul.
A mediana do Brasil como um todo foi 46,1, sendo a menor mediana encontrada na região Norte ( 41,8 ) e a maior encontrada na região Sudeste ( 47,0 ).
A nota mínima nesta parte foi zero em quase todas as regiões, exceto na região Centro-Oeste ( 7,5 ).
O Gráfico 3.5 propicia a avaliação do desempenho dos estudantes no Componente de Formação Geral a partir do histograma da distribuição das notas correspondentes.
A distribuição é unimodal, com moda em ( 40 ; 50 ], o mesmo intervalo modal da prova como um todo.
Nota -se, ainda, que no gráfico 3.5 as notas apresentam uma maior dispersão do que no Gráfico 3.1 ( distribuição das notas da prova ), confirmado pela comparação dos desvios padrões : 11,5 para a nota da prova como um todo e 14,4 para o Componente de Formação Geral.
Para o Componente de Formação Geral, o coeficiente de assimetria da distribuição das notas dos estudantes é negativo ( -0,14 ).
Por ser um coeficiente pequeno, a distribuição apresenta uma leve concentração à direita e cauda maior à esquerda, porém podemos considerar que a distribuição é aproximadamente simétrica.
Em todas as Grandes Regiões os histogramas também possuem assimetria levemente negativa ( entre -0,14 para a Sudeste e -0,01 para a Centro-Oeste ), com a região Sul com valor mais expressivo : -0,27.
Nos Gráficos 3.6, 3.7 e 3.8 são apresentadas as informações referentes ao desempenho dos Concluintes no Componente de Formação Geral, em diferentes agregações : Grande Região do país, Categoria Administrativa e Organização Acadêmica.
Observa -se pelo Gráfico 3.6 que existe diferença estatisticamente significativa ao nível de 95 \% entre a maior média das notas na Componente de Formação Geral, obtida na região Sudeste ( 46,7 ) e a menor, obtida na região Norte ( 41,7 ).
Vemos que o intervalo de confiança mais largo é o da região Norte ; já o intervalo mais estreito é observado na região Sudeste.
Este fato está relacionado, também, com o tamanho da população envolvida, menor na região Norte do que na Sudeste.
No Gráfico 3.7, que representa as notas médias no Componente de Formação Geral segundo Categoria Administrativa do país, observa -se que existe diferença estatisticamente significativa entre as médias.
Assim como ocorreu para as notas da prova como um todo, para o Componente de Formação Geral, os concluintes de Fisioterapia, das IES Públicas ( 50,7 ) obtiveram uma média maior do que os das IES Privadas ( 45,0 ).
Considerando -se o tipo de Organização Acadêmica, nota -se, no Gráfico 3.8, que existe uma diferença estatisticamente significativa entre a maior e a menor média.
Nas Universidades, a média ( 47,3 ) é mais elevada do que em Centros Universitários ( 45,3 ) e Faculdades ( 44,3 ).
A diferença entre esses dois últimos tipos de Organização Acadêmica é também estatisticamente significativa. 3.1.3 Estatísticas Básicas do Componente de Conhecimento Específico A Tabela 3.3 apresenta as estatísticas básicas referentes ao Componente de Conhecimento Específico da área de Fisioterapia.
A média do desempenho dos alunos do Brasil como um todo foi 49,4.
A maior média foi obtida na região Sul ( 52,9 ), e a menor, na região Norte ( 45,9 ).
As demais médias foram : 47,8 na região Nordeste, 49,9 na região Sudeste e 47,8 na região Centro-Oeste.
Quanto à variabilidade das notas, o desvio padrão do Brasil como um todo foi 13,1, sendo o maior desvio padrão observado na região Sudeste ( 13,2 ) e o menor na região Sul ( 12,3 ).
O desvio padrão das outras regiões foi : 12,6 na região Norte, 13,0 na região Nordeste e 12,9 na região Centro-Oeste.
A mediana das notas dos estudantes de todo o Brasil foi 49,6.
A maior mediana ocorreu na região Sul ( 53,3 ) e a menor na região Norte ( 46,1 ).
As demais medianas foram 47,9 na região Nordeste, 50,4 na região Sudeste e 47,4 na região Centro-Oeste.
A maior nota máxima do Brasil como um todo foi 91,8, sendo obtida por pelo menos um aluno da região Sul.
As demais notas máximas foram : 84,5 na região Norte, 90,6 na região Nordeste, 89,1 na região Sudeste e 85,9 na região Centro-Oeste.
A menor nota mínima foi zero nas regiões Norte e Nordeste, em quanto nas outras regiões as notas mínimas foram 7,4 na região Sudeste, 11,1 na região Sul e 4,2 na região Centro-Oeste.
Assim como os Gráficos 3.1 e 3.5, o Gráfico 3.9, apresentado a seguir, proporciona uma avaliação do desempenho de concluintes em relação ao Componente de Conhecimento Específico com um histograma da distribuição das notas correspondentes.
Esta também é uma distribuição unimodal, e o grupo modal é o ( 50 ; 60 ], acima do grupo modal para a Formação Geral.
O coeficiente de assimetria da distribuição das notas do Componente de Conhecimento Específico é negativo ( -0,12 ).
Nota -se pelo histograma ( Gráfico 3.9 ) que esta distribuição apresenta uma leve concentração à direita e cauda maior à esquerda.
Em quase todas as Grandes Regiões os histogramas também possuem assimetria negativa ( entre -0,18 para o Sudeste e -0,04 para o Norte ), exceto na região Centro-Oeste ( 0,03 ) que possui assimetria positiva, porém pequena.
Os Gráficos 3.10, 3.11 e 3.12 apresentam, respectivamente, uma comparação dos resultados em relação à Grande Região do país, à Categoria Administrativa e à Organização Acadêmica, agora levando em conta o desempenho de alunos no Componente de Conhecimento Específico da prova.
Pelo Gráfico 3.10, observa -se que existe diferença estatisticamente significativa entre a maior e a menor média das notas no Componente de Conhecimento Específico, observadas nas regiões Sul ( 52,9 ) e Norte ( 45,9 ) respectivamente.
Quanto à Categoria Administrativa ( Gráfico 3.11 ), observa -se um comportamento semelhante àquele da parte de Formação Geral e à prova como um todo, ou seja, existe diferença estatisticamente significativa entre as médias das IES Públicas e Privadas, sendo que a maior média foi obtida por alunos de IES Públicas de ensino ( 58,1 ).
Quanto ao Gráfico 3.12, observa -se, mais uma vez, que existe diferença estatisticamente significativa ao nível de 95 \% entre as notas no Componente de Conhecimento Específico dos diferentes tipos de Organização Acadêmica.
Sendo que a média dos concluintes das Universidades ( 51,8 ) foi maior do que de Centros Universitários ( 49,4 ) e de Faculdades ( 46,3 ). 3.2 ANÁLISE DAS QUESTÕES OBJETIVAS 3.2.1 Componente de Formação Geral A Tabela 3.4 apresenta as estatísticas básicas relativas às oito questões objetivas do componente da prova que abrange a Formação Geral dos estudantes.
A média do Brasil foi 47,4.
A menor média foi encontrada na região Norte ( 44,5 ) e a maior na região Sudeste ( 47,9 ).
As demais médias foram 47,3 nas regiões Nordeste e Centro-Oeste, e 47,6 na região Sul.
O desvio padrão do Brasil foi 17,4, sendo o maior desvio padrão encontrado nas regiões Norte, Nordeste e Centro-Oeste ( 17,5 ), e o menor na região Sul ( 17,2 ).
O desvio padrão na região Sudeste foi 17,3.
As mesmas medianas ( 50,0 ), notas máximas ( 100,0 ) e notas mínimas ( 0,0 ), nas questões objetivas do Componente de Formação Geral, foram obtidas em todas as regiões, sem exceção.
A Tabela 3.5 apresenta o índice de facilidade e o índice de discriminação ( ponto bisserial ) para cada uma das questões objetivas do Componente de Formação Geral.
Quanto ao índice de facilidade, foram usadas as seguintes cores para diferenciar o nível de dificuldade da questão : • Azul para as questões classificadas com índice muito fácil ( \&gt ; =0,86 ), verde para as questões classificadas com índice fácil ( 0,61 a 0,85 ), amarelo para as questões classificadas com médio ( 0,41 a 0,60 ), vermelho para as questões classificadas com difícil ( 0,16 a 0,40 ) e roxo para as questões classificadas com muito difícil ( \&lt ; =0,15 ).
Já quanto ao índice de discriminação, foram usadas as seguintes cores para qualificar a questão : • As questões classificadas com índice fraco receberam a cor vermelha ( \&lt ; =0,19 ), as classificadas com médio receberam a cor amarelo ( 0,20 a 0,29 ), as classificadas com bom receberam a cor verde ( 0,30 a 0,39 ) e as classificadas com muito bom ( \&gt ; =0,40 ) receberam a cor azul.
As questões objetivas do Componente de Formação Geral, segundo o índice de facilidade, foram assim avaliadas : das oito questões, uma teve o índice de facilidade classificado como muito fácil e outra como fácil.
Três questões foram tidas como médio, por terem índice de acertos situado na faixa entre 0,41 e 0,60 ( de 41,0 \% a 60,0 \% de acertos ).
Outras três questões foram consideradas de dificuldade difícil, situando -se no intervalo entre 0,16 e 0,40 do índice de facilidade, ou seja, houve entre 16,0 \% e 40,0 \% de acertos.
Como já comentado, para análise das questões objetivas relativas à Formação Geral segundo o poder de discriminação, utilizou -se, o índice de discriminação ponto bisserial.
Nesta análise, as questões foram assim avaliadas : cinco das oito questões apresentaram índices acima de 0,40 e, assim, foram classificadas com índice muito bom para esse grupo de alunos ; uma questão teve índice de discriminação bom, entre 0,30 e 0,39 ; e duas questões tiveram índice de discriminação médio, entre 0,20 e 0,29.
Nenhuma questão foi classificada com um índice fraco e, portanto, nenhuma questão foi eliminada pelo critério ponto bisserial.
O índice de facilidade variou de 0,18 a 0,96, e o de discriminação, de 0,23 a 0,48.
As questões com índice de discriminação muito bom, de números 3, 4, 6, 7 e 8, ficaram espalhadas entre os diferentes níveis de dificuldade : uma ( questão 8 ) classificada na categoria fácil, outras três na categoria médio do índice de facilidade e uma ( questão 7 ) na categoria difícil.
Em particular, a questão 3 apresentou maior poder discriminatório, com índice 0,48, e foi a quinta mais fácil, com uma proporção de 0,42 de acertos.
A questão de número 5 apresentou índice de facilidade 0,18, ou seja, um quantitativo de 18,0 \% de estudantes conseguiu resolvê -la, dentro do universo de participantes.
Além disso, seu índice de discriminação foi bom, com valor igual a 0,35.
O Gráfico 3.13, para exemplificar, analisa o comportamento da questão de número 3 de Formação Geral.
Trata -se da quinta questão mais fácil e a que obteve o maior índice de discriminação dessa parte da prova.
Neste gráfico, cada uma das cinco curvas representa o percentual de respostas em determinada alternativa da questão, em função da nota dos estudantes nesta parte da prova ( Formação Geral/Múltipla Escolha ), antes de possíveis eliminações pelo critério do ponto bisserial.
A curva em vermelho corresponde à alternativa E, a correta para esta questão.
Assim, observa -se que entre os estudantes com notas mais baixas, nessa parte do exame, a situação mais frequente foi a escolha de uma das alternativas incorretas : a alternativa A ( em azul ) ou C ( em preto ).
Na medida em que a nota aumenta, indicando desempenho melhor nesta parte da prova, aumenta concomitantemente a proporção de estudantes que selecionaram a alternativa correta E, atingindo 100 \% para os alunos com 8 acertos.
Essa análise permite verificar como a questão discriminou os grupos de desempenho, justificando o alto índice obtido na questão.
Os gráficos relativos às demais questões de Formação Geral constam do Anexo I. 3.2.2 Componente de Conhecimento Específico A Tabela 3.6 apresenta as estatísticas básicas em relação às questões objetivas do Componente de Conhecimento Específico da prova, por Grande Região.
A média do Brasil deste componente foi de 51,3.
A menor média foi observada na região Norte ( 47,2 ) e a maior na região Sul ( 54,8 ).
O desvio padrão de todo o Brasil foi 13,8, sendo o menor desvio padrão encontrado na região Sul ( 13,0 ) e o maior na região Sudeste ( 13,9 ).
A mediana de todo o Brasil foi 52,2.
Na região Sudeste a mediana também foi 52,2, mais alta ( 56,5 ) na região Sul e mais baixa ( 47,8 ) nas regiões Norte, Nordeste e Centro-Oeste.
A nota máxima da prova ( 95,7 ) foi obtida, nas questões objetivas do Componente de Conhecimento Específico, por pelo menos um aluno nas regiões Nordeste e Sul, ao passo que nas outras regiões as notas máximas foram 87,0 na região Norte, e 91,3 nas regiões Sudeste e Centro-Oeste.
A nota mínima foi zero nas regiões Norte e Nordeste, enquanto que nas outras regiões as notas mínimas foram 8,7 na região Sudeste, 13,0 na região Sul e 4,3 na região Centro-Oeste.
A Tabela 3.7 apresenta o índice de facilidade e o índice de discriminação ( ponto bisserial ) para cada uma das questões objetivas do Componente de Conhecimento Específico.
Para facilitar a diferenciação das questões usaram -se as mesmas cores da Tabela 3.5 para as diferentes classificações dos índices de facilidade e de discriminação.
Dentre as questões objetivas da parte da prova relativa ao Componente de Conhecimento Específico, nenhuma foi anulada pela Comissão.
Desse modo, a classificação quanto ao índice de facilidade foi estabelecida com base nas 27 questões.
A partir dos índices obtidos, pode -se concluir que uma boa parcela das 27 questões objetivas da prova foi considerada como difícil, dez ao todo.
As 17 questões restantes foram classificadas com índice de facilidade da seguinte maneira : duas questões muito fáceis, cinco fáceis, nove médias e uma muito difícil.
Já quanto aos índices de discriminação das questões objetivas do Componente de Conhecimento Específico da prova, tem -se como resultado a seguinte classificação : 11 das 27 questões foram consideradas como boas, enquanto nenhuma delas teve índice de discriminação muito bom.
Dentre as demais, 12 delas foram classificadas como médias e outras quatro como fracas, sendo 16, por conseguinte, a quantidade de questões nos dois patamares mais baixos de discriminação.
Constata -se, assim, que a prova – no que se refere ao Componente de Conhecimento Específico – possuía capacidade mediana de discriminar entre aqueles que dominam ou não o conteúdo.
Dentre as questões que alcançaram os maiores índices de discriminação, as de números 9, 14, 15, 18, 19, 24, 27, 28, 29, 31 e 35 classificadas com índice bom, situando -se no intervalo de 0,30 a 0,39 do índice, uma delas ( questão 15 ) foi classificada na categoria muito fácil, outra ( questão 35 ) como fácil, sete ( questões 9, 14, 18, 24, 27, 28 e 29 ) como média, duas ( questões 19 e 31 ) como difícil quanto ao índice de facilidade.
A questão de número 10 foi a mais difícil dentre as 27 questões específicas, com baixo índice de facilidade, apenas 14,0 \% de acertos.
Essa questão apresentou poder discriminatório igualmente baixo, 0,10, o que comprova ter sido esta a mais difícil para os estudantes.
Destacam -se, também, a questão 16, com índice de facilidade 0,17, o que, em termos percentuais, corresponde a 17,0 \% de estudantes que responderam acertadamente, obtendo, ainda, 0,08 de índice de discriminação, e a questão 20, com índice de facilidade 0,22, o que, em termos percentuais, corresponde a 22,0 \% de estudantes que responderam acertadamente, obtendo, ainda, 0,18 de índice de discriminação.
Tais questões foram, portanto, três das mais difíceis da prova.
As questões 10, 16 e 20 foram eliminadas do cômputo da nota final pelo critério do ponto bisserial.
Além destas três, outra questão com índice fraco de discriminação, questão 33 também não foi computada.
A título de exemplo das análises do comportamento das questões objetivas, o Gráfico 3.14 analisa a questão 24 do Componente de Conhecimento Específico.
Esta foi uma das questões classificadas como médio nesta prova, apresentando índice de facilidade 0,55, ou seja, 55,0 \% dos estudantes assinalaram acertadamente a opção A, correspondente ao gabarito.
Seu índice de discriminação foi igual a 0,39, classificado como bom.
Neste gráfico, cada uma das cinco curvas representa o percentual de respostas em determinada alternativa da questão 24, em função da nota dos estudantes nesta parte da prova, antes de possíveis eliminações de questões pelo critério do ponto bisserial.
A alternativa correta A, representada no gráfico pela curva em azul, foi escolhida em maiores proporções pelos alunos com desempenho melhor nesta parte da prova.
Já as alternativas incorretas, também denominadas distratores, foram selecionadas principalmente por aqueles com notas mais baixas.
A proporção de alunos que selecionou a resposta correta A aumenta gradativamente, chegando a atingir 100 \% para aqueles que acertaram pelo menos 23 questões, enquanto a proporção dos que escolheram alternativas incorretas decai, a partir da primeira nota não nula, como função da nota nesta parte da prova.
Os gráficos relativos às demais questões do Componente de Conhecimento Específico constam do Anexo I. 3.3 ANÁLISE DAS QUESTÕES DISCURSIVAS 3.3.1 Componente de Formação Geral As análises dos resultados de desempenho dos estudantes de Fisioterapia nas duas questões discursivas relativas à Formação Geral encontram -se na Tabela 3.8 e no Gráfico 3.15.
Na tabela 3.8 observa -se que as notas médias foram mais baixas nesse conjunto de questões do que no das objetivas.
Os estudantes de todo o Brasil obtiveram, em Formação Geral, média 47,4 nas questões objetivas e 43,3 nas questões discursivas.
A mediana também confirma o pior desempenho dos alunos de todo o Brasil nas questões discursivas do Componente de Formação Geral.
Enquanto essa estatística foi de 48,5 para questões discursivas, para as questões objetivas a mediana foi de 50,0.
Pode -se, também, notar um aumento do desvio padrão de 17,4, nas questões objetivas do Componente de Formação Geral dos alunos de todo o Brasil, para 20,9 nas questões discursivas do mesmo componente.
Como já comentado, a mediana de todo o Brasil, neste componente, foi 48,5, sendo a maior mediana encontrada na região Sudeste ( 50,0 ) e a menor mediana encontrada na região Norte ( 40,0 ).
A maior nota máxima foi obtida, nas questões discursivas do Componente de Formação Geral, por pelo menos um aluno da região Nordeste ( 92,0 ), enquanto a menor nota máxima foi obtida na região Norte ( 85,0 ).
Em todas as regiões a nota mínima foi zero, sem exceção.
O Gráfico 3.15 representa a distribuição das notas nas questões discursivas no Componente de Formação Geral.
A moda desta distribuição ocorre no intervalo ( 50 ; 60 ].
Devido à grande quantidade de notas zero e à alta frequência de alunos que deixaram este tipo de questão em branco o intervalo [ 0 ; 10 ] se apresenta como uma segunda moda.
A distribuição possui assimetria à esquerda, coeficiente de assimetria -0,76.
Em todas as regiões o coeficiente de assimetria também é negativo, o maior em módulo, -0,89, nas regiões Sudeste e Sul, e o menor na região Norte, -0,42.
A avaliação das questões discursivas de Formação Geral considerou em separado o conteúdo ( peso 0,8 na nota ) e o desempenho em Língua Portuguesa ( peso 0,2 ).
Na sequência, os resultados verificados para cada uma das questões discursivas de Formação Geral com relação ao conteúdo serão apresentados, estabelecendo -se relações com os temas abordados em cada uma delas.
Os comentários da Banca de docentes corretores a respeito do observado na correção das respostas dos estudantes, suas impressões e conclusões serão apresentados junto à análise de cada questão.
Em seguida será feita uma análise do desempenho de Língua Portuguesa.
Os comentários da Banca de docentes corretores serão apresentados para o conjunto de questões.
Cumpre esclarecer que, tendo em vista que as questões discursivas de Formação Geral são padronizadas, ou seja, constam de todas as provas, os comentários da Banca são os mesmos para todas as carreiras acadêmicas, sendo direcionados a todos os estudantes que participaram do ENADE/2013.
A seguir, serão analisados os desempenhos da Área de Fisioterapia nas duas questões discursivas de Formação Geral do ENADE/2013, comparando os resultados obtidos com comentários para cada questão. 3.3.1.1 Análise de Conteúdo Questão Discursiva 1 do Componente de Formação Geral Os dados de Fisioterapia, obtidos a partir das respostas à questão 1, encontram -se na Tabela 3.9 e no Gráfico 3.16.
Nessa questão – de melhor desempenho dentre as duas de Formação Geral – os alunos de todo Brasil tiveram média, 45,4.
A maior média para a questão 1 foi obtida na região Sudeste ( 46,2 ), e a menor, na região Norte ( 41,7 ).
Quanto à variabilidade das notas, o desvio padrão de todo o Brasil foi 25,4.
O menor desvio padrão foi obtido na região Sul ( 24,0 ) e o maior desvio padrão foi obtido na região Norte ( 27,0 ).
A mediana da questão discursiva 1 foi 50,0 em quase todas as regiões do Brasil, exceto na região Norte ( 45,0 ).
As notas máximas ( 100,0 ) foram as mesmas para todas as regiões do Brasil, sem exceção.
Além disso, a nota mínima foi zero em todas as regiões do país, sem exceção.
O Gráfico 3.16 mostra a distribuição das notas na questão discursiva 1 do Componente de Formação Geral.
Observa -se que a maior frequência corresponde aos alunos obtiveram nota no intervalo ( 40 ; 50 ], correspondendo à moda da distribuição.
Destacam -se os alunos que deixaram a questão em branco e o intervalo ( 60 ; 70 ] como máximos locais.
A distribuição das notas possui assimetria negativa ( -0,57 ), assim como em todas as regiões do Brasil ( entre -0,60 para o Sudeste e -0,20 para o Norte ). 3.3.1.2 Comentários sobre a correção de Conteúdo das respostas à Questão Discursiva 1 O enunciado restou claro, elaborado com linguagem técnica, estando inserido nas Diretrizes gerais, com nível médio de profundidade, já que abordava tema vinculado à Área de formação da maior parte dos examinandos.
A divisão das perguntas também ficou clara indicando, passo a passo, o caminho de resposta a seguir, com base no texto.
Verificou -se dificuldade quanto à análise dos programas de governo existentes que contribuem para melhoria do saneamento básico.
Nesse ponto, os alunos referiram, com muita frequência, a ações relacionadas à Saúde, tais como : médico da família, ações educativas, ações ambientais, melhoria nas condições de habitação ( sem necessariamente vincular a programas de governo federal ou estaduais ou municipais ).
Os programas citados foram : Programa Saúde da Família ( PSF ), Programa de Aceleração do Crescimento ( PAC ) e Minha Casa Minha Vida e, ainda, o SUS.
Apesar de uma minoria de examinandos ter indicado o PAC do Governo Federal, a maior parte deles apresentou alguma solução para a questão do saneamento básico.
Os estudantes indicaram, com facilidade, a questão da necessidade do tratamento da água como item essencial do saneamento básico.
Muitos concluintes, através da sua experiência local, citaram situações em que as crianças caminhavam de pés no chão perto de rotas de esgoto a céu aberto.
Também criticaram a inércia governamental em não combater as invasões ilícitas, em locais inadequados, sem estrutura de água potável e também de saneamento básico.
Diante das observações preliminares pode ser extraída uma classificação das respostas obtidas que ficou estabelecida em fracas, médias e boas.
As notas consideradas fracas ( notas de 0 a 3,0 ) foram conferidas àqueles que não conseguiram relacionar o saneamento básico com a importância do tratamento de água, os projetos para resolver o problema e nem indicaram soluções para a questão.
As respostas que tiveram média avaliação ( notas de 3,5 a 7,0 ) tocaram, de algum modo, nos pontos indicados, mesmo que de forma insuficiente.
Já as respostas consideradas boas ( notas de 7,5 a 10,0 ) preencheram os requisitos do padrão de resposta, não logrando nota máxima, na maior parte dos casos, por ausência de indicação de programa governamental, mas sugerindo soluções técnicas para o problema.
As respostas mais comuns foram as consideradas medianas e que indicaram, com frequência, que o problema do saneamento básico estaria nas classes de renda mais desfavorecida e nas regiões consideradas como periferias das grandes cidades.
Também foram frequentes respostas que citaram a necessidade de tratamento de água, acompanhada da necessidade de maiores investimentos nessa área. condições de habitação ( sem necessariamente vincular a programas de governo federal ou estaduais ou municipais ).
Os programas citados foram : Programa Saúde da Família ( PSF ), Programa de Aceleração do Crescimento ( PAC ) e Minha Casa Minha Vida e, ainda, o SUS.
Apesar de uma minoria de examinandos ter indicado o PAC do Governo Federal, a maior parte deles apresentou alguma solução para a questão do saneamento básico.
Os estudantes indicaram, com facilidade, a questão da necessidade do tratamento da água como item essencial do saneamento básico.
Muitos concluintes, através da sua experiência local, citaram situações em que as crianças caminhavam de pés no chão perto de rotas de esgoto a céu aberto.
Também criticaram a inércia governamental em não combater as invasões ilícitas, em locais inadequados, sem estrutura de água potável e também de saneamento básico.
Diante das observações preliminares pode ser extraída uma classificação das respostas obtidas que ficou estabelecida em fracas, médias e boas.
As notas consideradas fracas ( notas de 0 a 3,0 ) foram conferidas àqueles que não conseguiram relacionar o saneamento básico com a importância do tratamento de água, os projetos para resolver o problema e nem indicaram soluções para a questão.
As respostas que tiveram média avaliação ( notas de 3,5 a 7,0 ) tocaram, de algum modo, nos pontos indicados, mesmo que de forma insuficiente.
Já as respostas consideradas boas ( notas de 7,5 a 10,0 ) preencheram os requisitos do padrão de resposta, não logrando nota máxima, na maior parte dos casos, por ausência de indicação de programa governamental, mas sugerindo soluções técnicas para o problema.
As respostas mais comuns foram as consideradas medianas e que indicaram, com frequência, que o problema do saneamento básico estaria nas classes de renda mais desfavorecida e nas regiões consideradas como periferias das grandes cidades.
Também foram frequentes respostas que citaram a necessidade de tratamento de água, acompanhada da necessidade de maiores investimentos nessa área. 3.3.1.3 Análise de Conteúdo da Questão Discursiva 2 do Componente de Formação Geral A Tabela 3.10 mostra que o desempenho dos estudantes na questão 2 ( média 39,1 ) foi inferior ao obtido na questão de número 1 ( média 45,4 ).
A região Sudeste foi aquela onde a média, nessa questão, foi maior ( 41,9 ), e a de menor média foi a região Norte ( 28,3 ).
Quanto à variabilidade das notas, o desvio padrão de todo o Brasil foi 26,0, superior ao obtido na questão de número 1 ( 25,4 ).
O maior desvio nessa questão foi obtido na região Nordeste ( 26,8 ), enquanto o menor foi obtido na região Centro-Oeste ( 24,6 ).
A mediana foi 45,0 para o Brasil como um todo, sendo encontrada, também, na região Centro-Oeste.
As medianas nas demais regiões foram : 25,0 na região Norte, 40,0 na região Nordeste e 50,0 nas regiões Sudeste e Sul.
As notas máximas ( 100,0 ) foram as mesmas em todas as regiões do Brasil, assim como as notas mínimas ( 0,0 ), sem exceção.
O Gráfico 3.17 mostra a distribuição das notas na questão discursiva 2 do Componente de Formação Geral.
Neste gráfico, observa -se a grande quantidade de alunos que deixaram a questão 2 em branco ( em torno de 18 \% ).
A moda da distribuição ficou no intervalo ( 40 ; 50 ], com um percentual um pouco superior a 18 \%.
O coeficiente de assimetria das notas nesta questão é negativo ( -0,18 ), indicando que a distribuição é assimétrica à esquerda.
O mesmo ocorre para quase todas as regiões da Brasil ( -0,33 na região Sudeste e -0,30 nas regiões Sul e Centro-Oeste ), exceto as regiões Norte ( 0,38 ) que possui assimetria positiva e Nordeste ( 0,01 ), com assimetria virtualmente nula. 3.3.1.4 Comentários sobre a correção de Conteúdo das respostas à Questão Discursiva 2 O tema abordado no enunciado, apesar da complexidade, é extremamente atual.
Essa atualidade facilitou a elaboração de alguma resposta, permitindo que todos os que se dispuseram a responder pudessem apresentar, pelo menos, alguns dos elementos integrantes do padrão de resposta.
A correção, seguindo o padrão de resposta, avaliou a capacidade de elaborar um texto dissertativo que abordasse o tema sob três pontos de vista : o papel da tecnologia digital na atualidade ; a garantia dos direitos do cidadão e do Estado diante do avanço tecnológico ; e o problema da segurança digital para o indivíduo, para o Estado e para as organizações ( empresas e/ou instituições nacionais ou internacionais ).
Dada a amplitude de respostas possíveis, os alunos tiveram desempenho satisfatório.
Identificamos que os concluintes, basicamente, trilharam dois caminhos : a ) a análise da influência das inovações tecnológicas, notadamente através das redes sociais, numa perspectiva essencialmente privada, daí defluindo as questões da privacidade e da segurança ; b ) uma perspectiva de segurança nacional-estatal com criticas à atuação americana de invasão de privacidade dos dados governamentais, e à inercia do governo brasileiro em reagir ao ataque, considerado, pela maioria esmagadora dos concluintes, como indevido e antiético.
Poucos foram os que uniram as duas facetas da questão : o público e o privado.
Na maior parte das provas, o concluinte indicou a ausência de segurança e de privacidade decorrentes das próprias atividades individuais.
Muitos respondentes citaram que, de forma voluntária, os próprios cidadãos enfraquecem a proteção à privacidade, pois divulgam sua vida pessoal, aspectos de seu cotidiano, por meio de imagens e do fornecimento de dados pessoais.
No aspecto segurança, foram inúmeras as indicações de acessos indevidos a dados pessoais com a perpetração de fraudes de todos os tipos, notadamente as bancárias e comerciais, com a aquisição de bens à revelia dos titulares dos dados.
Também no tema segurança, muitos indicaram o aspecto positivo decorrente da vigilância generalizada, não somente através da rede de computadores, mas também pela quantidade de câmeras existentes, que contribuem para a identificação de criminosos procurados ou dos que realizam crimes pela internet.
Na perspectiva pública, a esmagadora maioria indicou repulsa pelo comportamento americano de invasão dos dados do governo brasileiro, reclamando reação do mesmo.
Alguns poucos justificaram a intervenção americana pela “ natural ” disputa de poder pelos Estados imperialistas. 3.3.1.5 Análise de Língua Portuguesa das Questões Discursivas do Componente de Formação Geral Os dados de Fisioterapia, obtidos a partir das respostas às questões discursivas do Componente de Formação Geral no que tange à Língua Portuguesa, encontram -se na Tabela 3.11 e no Gráfico 3.18.
Nesse aspecto, os alunos de todo Brasil tiveram média, 47,4.
A maior média com respeito à Língua Portuguesa foi obtida na região Sul ( 49,8 ), e a menor, na região Nordeste ( 44,9 ).
Quanto à variabilidade das notas, o desvio padrão de todo o Brasil foi 21,6.
O menor desvio padrão foi obtido na região Centro-Oeste ( 20,1 ) e o maior desvio padrão foi obtido na região Nordeste ( 23,0 ).
A mediana da nota de Língua Portuguesa foi 52,5 para as regiões Norte, Nordeste e Centro-Oeste, mais alta para as regiões Sudeste e Sul ( 55,0 ).
Para o Brasil como um todo foi também 52,5.
A maior nota máxima foi obtida por pelo menos um aluno da região Nordeste ( 92,5 ), enquanto que nas outras regiões as notas máximas foram : 82,5 na região Norte, 90,0 na região Sudeste, e 85,0 nas regiões Sul e Centro-Oeste.
Além disso, a nota mínima foi zero em todas as regiões do país, sem exceção.
Norte, 90,0 na região Sudeste, e 85,0 nas regiões Sul e Centro-Oeste.
Além disso, a nota mínima foi zero em todas as regiões do país, sem exceção.
O Gráfico 3.18 mostra a distribuição das notas de Língua Portuguesa do Componente de Formação Geral.
Observa -se que a maior frequência corresponde aos alunos que obtiveram nota no intervalo ( 50 ; 60 ], correspondendo à moda da distribuição.
A distribuição das notas possui assimetria negativa ( -1,02 ), assim como em todas as regiões do Brasil ( entre -1,20 para o Sul e -0,77 para o Nordeste ). 3.3.1.6 Comentários sobre a correção das respostas de Formação Geral com respeito à Língua Portuguesa Um aspecto interessante a ser destacado é a diferença de desempenho entre a questão 1 e a questão 2, principalmente nos aspectos textuais – os textos produzidos para a questão 1 apresentaram resultados melhores do que os produzidos para a questão 2.
Uma hipótese que provavelmente explique essa discrepância é o fato de o tema “ saneamento básico ” ser mais relevante para profissionais de saúde do que o tema “ segurança e privacidade ”, constituindo tema trabalhado no currículo escolar.
Prova disso é que quase todos os textos da questão 1 citam termos específicos da área de saúde, como “ contaminação feco-oral ”, “ higiene corporal ”, “ ingestão de alimentos ou água contaminada por vetores biológicos ”.
Os resultados da avaliação correspondem aos seguintes aspectos observados em cada competência : a ) aspectos ortográficos : o desempenho dos participantes ficou concentrado, majoritariamente, nos níveis 3 ( até 3 desvios ) e 2 ( de 4 a 7 desvios ), devido ao grande índice de desvios de acentuação, contrabalançada pelo pequeno número de desvios em grafia.
Foi reduzido o número de textos com ausência de desvios ( enquadrados no nível 4 ).
Alguns textos, com pior desempenho, foram enquadrados no nível 1 ( de 8 a 12 desvios ), enquanto o nível zero foi atribuído a pouquíssimos casos, já que ele revela falta absoluta de domínio das convenções ortográficas.
Observou -se, portanto, que existe uma diferença muito grande de desempenho nos dois aspectos analisados : baixíssimo índice de desvios ortográficos e grande índice de desvios de acentuação.
Em vários casos, ocorre ausência completa de acentuação gráfica.
Os resultados revelam, portanto, que a tendência dominante entre os universitários brasileiros é a eliminação da acentuação gráfica, talvez motivada pela ausência de esclarecimento dos meios de comunicação, das autoridades e das escolas sobre as decisões do Acordo Ortográfico da Língua Portuguesa de 1990.
Os casos mais sistemáticos de ausência de acentuação são :  palavras proparoxítonas ( “ politicas ”, “ publicas ”, “ juridicas ” ) ;  palavras paroxítonas terminadas em ditongo crescente ( “ varios ”, “ Brasilia ”, “ familia ” ) ;  palavras oxítonas ( “ ninguem ”, “ esta ”, “ ate ” ).
Quanto ao domínio das convenções relativas à grafia das palavras, são pontuais os problemas encontrados ( “ infecciozas ”, “ obtos ”, “ hipedemia ”, “ saniamento ”, “ cituação ”, “ abitantes ” ).
Os únicos desvios recorrentes foram : a ausência de segmentação nos casos de verbo + pronome oblíquo “ se ”, combinada, eventualmente, à alteração da grafia do pronome ( “ percebese ”, “ percebece ”, “ preferece ” ) ; a grafia da forma verbal de pretérito imperfeito do subjuntivo como se fosse uma ocorrência de verbo + pronome oblíquo ( “ distribui -se ” por “ distribuísse ” ) Vale observar, também, que, ao contrário do que se esperava, não apareceram abreviaturas próprias do “ internetês ”, ou seja, dos hábitos de comunicação escrita adquiridos pelo uso de redes sociais e emails. b ) aspectos textuais : o desempenho dos participantes ficou concentrado, majoritariamente, nos níveis 3 e 2, devido à grande ocorrência de problemas de estruturação textual.
Foi muito baixo o número de textos enquadrados no nível 4 por não apresentar qualquer problema estrutural.
Alguns textos com pior desempenho foram enquadrados no nível 1, enquanto o nível zero foi atribuído a pouquíssimos casos, já que ele revela a existência de um texto sem articulação e com comprometimento do sentido.
Observou -se que a grande maioria dos participantes não distribuiu as ideias em parágrafos, talvez devido ao pequeno número de linhas disponibilizadas para a resposta da questão ou, quem sabe, pela suposição de que não seria necessária essa divisão por não se tratar de um texto no modelo de uma redação dissertativo-argumentativa, como solicitado nos vestibulares.
Essa competência é a que se revela como a mais problemática entre os participantes, porque são muitos os problemas observados, desvios acumulados durante toda a formação do estudante e que não se resolvem com um estudo autodidata, como acontece com regras ortográficas ou morfossintáticas : sequência justaposta de ideias sem encaixamentos sintáticos ; redução drástica de estruturas subordinadas, ao lado do aumento na frequência de estruturas coordenadas e absolutas ; redução no uso de conectores para expressar relações lógicas essenciais à construção do texto, substituídas pela exigência de inferência por parte do interlocutor para suprir a sua ausência ; emprego equivocado de operadores que não estabelecem relações lógicas coerentes entre ideias do texto ; emprego inadequado do pronome relativo ( com omissão da preposição ou a utilização de pronome inadequado, como “ onde ” ) ; repetição exaustiva de termos sem a utilização de procedimentos mais sofisticados de substituição ( hiperonímias, hiponímias, nominalizações, expressões metafóricas ) ; frases fragmentadas que comprometem a estrutura lógico-gramatical ; frases formadas apenas por oração subordinada, sem oração principal.
Um importante aspecto a destacar é o baixíssimo desempenho de uma parte dos participantes em relação à estrutura formal do texto produzido, o que é extremamente preocupante ao se levar em conta que são graduandos em fase final de formação.
São frequentes os casos de desvios de estruturação frasal, com uso inadequado ou ausência de conectivos entre parágrafos e entre frases.
Em uma parte dos textos, falta um mínimo de textualidade e de domínio do registro padrão da língua.
Na verdade, observam -se relações linguísticas quase agramaticais, como as estabelecidas pela sequência de gerúndios sem o apoio de um ponto de partida para a organização das informações gramaticais e semânticas.
Quanto à utilização dos mecanismos de referenciação, dois aspectos devem ser destacados : a ocorrência, em uma boa parte dos textos, de repetições de palavras ou expressões sem a utilização de termos sinônimos ou pronomes, como seria adequado ; e a alternância entre os pronomes demonstrativos este / esse, e suas formas flexionadas, sem a observância de um critério de utilização.
Esse último aspecto tem se generalizado na modalidade oral da língua e, cada vez mais, tem se tornado uma constante na modalidade escrita de jornais e revistas : predominância da forma “ esse ” na modalidade oral e da forma “ este ” na modalidade escrita, provavelmente por um processo de hipercorreção. c ) aspectos morfossintáticos e vocabulares : o desempenho dos participantes ficou concentrado, majoritariamente, nos níveis 3 e 2, devido à grande ocorrência de problemas de regência, concordância e uso adequado dos sinais de pontuação.
O nível 4 foi atribuído a um número muito reduzido de textos.
Alguns textos com pior desempenho foram enquadrados no nível 1, enquanto o nível zero foi atribuído a pouquíssimos casos, já que ele revela a existência de um texto sem articulação e com comprometimento do sentido.
Os resultados são muito transparentes em relação aos aspectos mais problemáticos no desempenho dos participantes.
O desvio mais frequente é a falta do sinal indicativo da crase, o que se configura como um problema de regência, verbal ou nominal, conforme o caso – isso revela que o usuário não tem consciência de que, sob a forma do termo “ a ” existe a presença de uma preposição “ a ”, exigida pela regência do termo anterior.
Embora em outros exames, como o Enem, a falta de crase seja apenada em acentuação, nesta correção amostral esse desvio foi considerado no âmbito dos aspectos morfossintáticos.
Outro problema relacionado à regência verbal e à nominal, encontrado frequentemente nas questões, foi a ausência de preposição antes de pronome relativo, processo generalizado na modalidade oral da língua, em situações de registro informal.
Apesar da possibilidade de que essa alteração de regência se generalize no padrão escrito da Língua Portuguesa, como já está ocorrendo até em textos jornalísticos, o não emprego da preposição foi apenado neste processo de avaliação.
Outro desvio muito frequente diz respeito aos processos de concordância verbal e de concordância nominal.
Quanto à concordância de número, observou -se ausência de marca ( com sujeito anteposto ou posposto ) ou uso indevido ( uso inadequado da marca de plural comandado pelo núcleo plural da locução adjetiva, apesar de o substantivo que funciona como núcleo do sintagma nominal estar no singular ).
Uma ocorrência que se destacou foi a ausência de acento circunflexo na forma plural do presente do indicativo do verbo “ ter ”, que foi considerada como um desvio na concordância verbal e não na acentuação gráfica.
Quanto à concordância de gênero, vários casos foram observados, normalmente no âmbito de sintagmas nominais longos, em que o adjetivo está afastado do substantivo.
Quanto à questão da colocação pronominal, foram poucos os casos observados.
Apesar de serem aspectos relacionados à oralidade, concluiu -se que, no registro escrito formal, a maioria dos participantes já incorporou regras como a não introdução da frase por um pronome oblíquo e a próclise na presença de um termo atrator.
Não se adotou, entretanto, o padrão excessivamente formal descrito pelas gramáticas normativas em relação à posição do pronome oblíquo em locuções verbais, já que esse uso está muito distante da prática cotidiana, até em textos mais formais.
Quanto aos aspectos vocabulares, alguns tipos de inadequação foram observados : expressões da oralidade apareceram em algumas respostas, mas sem maior relevância do ponto de vista quantitativo ; seleção vocabular incompatível com o contexto, gerando situações de falta de inteligibilidade ; falta de domínio de vocabulário mais abstrato e de maior complexidade, essencial ao desenvolvimento do texto de base dissertativa.
Quanto aos aspectos vocabulares, alguns tipos de inadequação foram observados : expressões da oralidade apareceram em algumas respostas, mas sem maior relevância do ponto de vista quantitativo ; seleção vocabular incompatível com o contexto, gerando situações de falta de inteligibilidade ; falta de domínio de vocabulário mais abstrato e de maior complexidade, essencial ao desenvolvimento do texto de base dissertativa.
Quanto à utilização dos sinais de pontuação, observou -se uma grande precariedade nos textos analisados.
Não foi apenada a ausência de vírgula para destacar locuções ou adjuntos adverbiais de pequena extensão deslocados de posição na frase, por ser um uso opcional.
São os seguintes os tipos de problemas encontrados : a ) vírgula : utilização de vírgula para separar o sujeito e o predicado ; ocorrência de apenas uma das vírgulas para separar uma palavra, uma expressão ou uma oração encaixada ; uso de vírgula no lugar do ponto para separar ideias que constituem períodos distintos ; ausência de vírgula para separar enumerações ; ausência de vírgula para separar oração adjetiva explicativa ou utilização inadequada para separar oração adjetiva restritiva ; b ) ponto e vírgula : utilização do ponto e vírgula no lugar de vírgula ; c ) ponto final : ausência de ponto final para separar períodos.
Várias marcas de oralidade foram identificadas, embora não em alta frequência : o uso do pronome relativo “ onde ” como relativo universal, falta de artigo definido antes de substantivo, repetição de palavras por falta de vocabulário, reduções como “ tá ”, “ pra ”, “ pro ”, “ prum ”, expressões informais como “ estão nem aí ”, “ tem que ”, “ levando com a barriga ”, “ não tá nem aí ”, “ a coisa fica séria ”, “ e aí haja dor de cabeça ”. 3.3.2 Componente de Conhecimento Específico Na parte da prova relativa às questões discursivas no Componente de Conhecimento Específico ( Tabela 3.12 ), observa -se que a média foi mais baixa do que nas questões discursivas do Componente de Formação Geral.
Enquanto no Componente de Formação Geral a média para estudantes de Fisioterapia de todo o Brasil foi 43,3, na parte de Conhecimento Específico a média foi 38,3.
A maior média deste componente foi obtida pelos estudantes da região Sul ( 42,2 ), e a menor, pelos da região Nordeste ( 36,1 ).
Quanto à variabilidade das notas, o desvio padrão de todo o Brasil foi 18,6.
O maior desvio padrão foi encontrado nas regiões Norte e Nordeste ( 19,2 ), e o menor, na região Sul ( 17,7 ).
A maior nota máxima, nas questões discursivas do Componente de Conhecimento Específico, foi obtida na região Nordeste ( 96,7 ), enquanto a menor nota máxima foi encontrada na região Centro-Oeste ( 80,0 ).
Além disso, a nota mínima ( 0,0 ) foi obtida por alunos de todas as regiões do Brasil sem exceção.
A mediana do Brasil como um todo foi 40,0, também encontrada nas regiões Norte, Sudeste e Centro-Oeste.
As demais medianas foram : 36,7 na região Nordeste e 43,3 na região Sul.
O Gráfico 3.19 representa a distribuição das notas nas questões discursivas no Componente de Conhecimento Específico.
Essa distribuição tem moda no intervalo ( 40 ; 50 ], mas o intervalo contíguo anterior, ( 30 ; 40 ], apresenta valor bem parecido.
O coeficiente de assimetria é negativo para todo o Brasil ( -0,22 ) e para todas as regiões ( entre -0,32 para a região Sul e -0,10 para a região Nordeste ), indicando uma assimetria à esquerda. 3.3.2.1 Análise da Questão Discursiva 3 do Componente de Conhecimento Específico Na questão discursiva 3, cujos resultados aferidos encontram -se descritos na Tabela 3.13, a média dos estudantes de todo o Brasil foi 40,4.
A menor média nessa questão foi obtida pelos alunos da região Nordeste ( 38,6 ), enquanto a maior média foi obtida na região Norte ( 43,4 ).
Quanto à variabilidade das notas, o desvio padrão de todo o Brasil foi 23,9.
O maior desvio padrão foi obtido na região Norte ( 27,8 ), enquanto o menor foi obtido na região Sul ( 22,6 ).
A maior nota máxima foi alcançada por pelo menos um aluno nas regiões Norte e Nordeste ( 100,0 ), enquanto nas outras regiões as notas máximas foram 95,0 na região Sudeste, e 90,0 nas regiões Sul e Centro-Oeste.
A mediana do Brasil como um todo foi 40,0, sendo atingida nas regiões Nordeste e Sudeste, enquanto que nas regiões Norte, Sul e Centro-Oeste, a mediana foi 45,0.
A nota mínima do Brasil como um todo foi zero, o mesmo valor para todas as regiões, sem exceção.
O Gráfico 3.20 mostra a distribuição das notas na questão discursiva 3, do Componente de Conhecimento Específico.
Essa distribuição apresenta uma moda principal no intervalo ( 50 ; 60 ], mas modas secundárias no intervalo ( 30 ; 40 ] e em questões em branco.
O coeficiente de assimetria é negativo para o Brasil ( -0,15 ) e para quase todas as regiões ( entre -0,30 na região Sul e -0,05 na região Nordeste ), exceto a região Norte ( 0,00 ) que possui assimetria nula. 3.3.2.2 Comentários sobre as respostas à Questão Discursiva 3 O enunciado da Questão 3 era claro e objetivo, sem deixar margem para duplas interpretações.
O contexto apresentava situações que levariam o aluno à resposta por meio de uma leitura cuidadosa apenas.
Sua objetividade permitiria respostas claras e concisas.
A única ressalva foi o uso da expressão “ 1º PO ”, que poderia ter sido apresentada de forma mais explícita “ 1ºdia de PO ”, para eliminar qualquer tipo de dúvida.
O padrão de resposta era claro e objetivo.
Cabe destacar que todos os itens da resposta para item ‘ a ’ estavam descritos no próprio texto, bastava o aluno relacionar ao que leu.
Em relação ao item ‘ b ’, por possibilitar várias técnicas como resposta, o padrão admitia uma ampla gama de opções.
A flexibilidade do padrão foi confirmada, pois não houve respostas corretas propostas de condutas fisioterapêuticas diferentes das previstas no padrão.
No entanto, alguns alunos usaram termos não reconhecidos e técnicas, batizadas com nomes de autores, também não referenciadas no padrão de respostas. 3.3.2.3 Análise da Questão Discursiva 4 do Componente de Conhecimento Específico A Tabela 3.14 contém as informações relativas à questão discursiva 4 do conjunto de questões do Componente de Conhecimento Específico.
O desempenho dos estudantes de todo o Brasil nesta questão foi inferior ao desempenho na questão discursiva 3.
A média geral do Brasil foi 33,1, sendo a menor média registrada na região Norte ( 28,5 ) e a maior na região Sul ( 38,2 ).
Quanto à variabilidade das notas, o desvio padrão de todo o Brasil foi 25,1.
O maior desvio padrão foi obtido na região Nordeste ( 25,8 ), enquanto o menor foi obtido na região Centro-Oeste ( 24,0 ).
A nota máxima ( 100,0 ) foi atingida em todas as regiões, sem exceção.
A mediana foi 40,0 para o Brasil com um todo, sendo encontrada também nas regiões Sudeste e Sul, enquanto nas outras regiões as medianas foram : 30,0 nas regiões Norte e Nordeste, e 35,0 na região Centro-Oeste.
A nota mínima foi zero para todas as regiões do Brasil, sem exceção.
O Gráfico 3.21 representa a distribuição de notas na questão discursiva 4, no Componente de Conhecimento Específico.
Essa distribuição é multimodal com moda principal no intervalo ( 30 ; 40 ].
O coeficiente de assimetria é positivo para todo o Brasil ( 0,29 ) e para todas as regiões ( entre 0,09 para a região Sul e 0,55 para a região Norte ), indicando uma assimetria à direita. 3.3.2.4 Comentários sobre as respostas à Questão Discursiva 4 A Questão 4 era clara e objetiva.
Solicitava ao aluno a identificação rápida do quadro para que respondesse adequadamente à questão.
A questão avaliava um nível elementar de conhecimento que todo fisioterapeuta deve possuir, envolvendo vivência prática.
A única dificuldade pode ter sido a necessidade do uso de nomenclaturas exigidas como resposta correta pelo padrão estabelecido.
Em muitas respostas verificou -se que houve desatenção ao enunciado quanto à classificação do tipo de hipertonia.
Além disso, a maioria dos estudantes não utilizou os termos técnicos corretos.
Em relação ao item ‘ b ’, a forma de identificação da hipertonia foi indicada por muitos alunos por escalas de avaliação de tônus, o que não foi considerado por não constar do padrão de respostas aprovado e divulgado.
Foi muito frequente o uso de termos não científicos, tanto ao avaliar o tipo de hipertonia quanto na descrição das técnicas.
A identificação da Escala de Arsworth, como critério de avaliação do tônus foi bastante frequente, e apesar de entendermos que essa escala quantifica e não avalia, consideramos que tal resposta poderia ser parcialmente 3.3.2.5 Análise da Questão Discursiva 5 do Componente de Conhecimento Específico A Tabela 3.15 contém as informações relativas à questão discursiva 5 do conjunto do Componente de Conhecimento Específico.
O desempenho dos estudantes nessa questão foi superior ao das questões discursivas 3 e 4.
A nota média dos estudantes de todo o Brasil foi 41,5.
A maior média foi registrada na região Sul ( 45,9 ), enquanto a menor média foi registrada na região Nordeste ( 39,5 ).
Quanto à variabilidade das notas, o desvio padrão dos alunos do Brasil, como um todo, foi 28,1, também encontrada na Sul.
Enquanto o maior desvio foi encontrado na região Norte ( 28,9 ), o menor foi encontrado na região Sudeste ( 27,6 ).
A nota máxima ( 100,0 ) foi alcançada em todas as regiões do Brasil, sem exceção.
A mediana foi 40,0 para o Brasil e também para todas as regiões, sem exceção.
A nota mínima foi zero para todas as regiões do Brasil, sem exceção.
O Gráfico 3.22 representa a distribuição das notas da questão discursiva 5 do Componente de Conhecimento Específico.
Essa distribuição apresenta várias modas, mas com moda principal no intervalo ( 30 ; 40 ].
Entre as três questões discursivas do Componente de Conhecimento Específico, esta foi a que apresentou um maior espalhamento das notas daqueles que resolveram a questão.
Além da moda principal, os máximos locais estavam localizados em questões em branco ( acrescidos de nota zero ) e ( 70 ; 80 ].
O coeficiente de assimetria para todos os alunos é 0,13 e para quase todas as regiões se mantém positivo, variando entre 0,04 nas regiões Norte e Centro-Oeste e 0,23 na região Nordeste, exceto na região Sul ( -0,02 ) que possui coeficiente de assimetria negativo, porém pequeno. 3.3.2.6 Comentários sobre as respostas à Questão Discursiva 5 A Questão 5 era clara e objetiva na descrição do contexto de um paciente.
O que a questão solicitava dos alunos nos itens ‘ a ’ e ‘ b ’ também não dava margem a interpretações diferentes do que o padrão de resposta apresentava.
No item ‘ c ’, pelo texto apresentado, observou -se que muitos alunos entenderam que deveriam apresentar um critério para identificação dos dois tipos de exercícios e, não, um critério para cada exercício como o padrão de resposta previa.
A Questão 5, além de ser a última da prova, exigia também resposta muito longa.
Possivelmente por estas razões a maioria dos alunos parece ter “ desistido ” de responder a tudo que era solicitado.
Acreditamos que muitos poderiam desenvolver mais suas respostas e evidenciar conhecimento sobre o tema se fossem feitos questionamentos mais específicos.
O padrão de respostas era flexível e abrangente ao apresentar, para o item ‘ a ’, sete opções de recursos poderiam constar das respostas dos alunos, e para o item ‘ b ’, cinco opções de respostas com as respectivas explicações, dentre as quais duas opções de respostas poderiam ser consideradas.
No entanto, infelizmente, como se exigiam muitos aspectos, os textos dos estudantes que chegaram a responder ficaram muito aquém do esperado.
O que mais se encontrou foram respostas com utilização de vários recursos indevidos no item ‘ a ’, explicações insuficientes no ‘ b ’ e respostas não precisas no item ‘ c ’.
Quanto ao item ‘ c ’, consideramos que houve falta de clareza e objetividade na formulação e, por isso, foi preciso aceitar uma gama muito grande de respostas como corretas.

% matched lemmas: 
\end{textsample}
